\documentclass[10pt]{article}
\usepackage[T1]{fontenc}
\usepackage[utf8]{inputenc}
\usepackage[french]{babel}
\usepackage{amsmath,amssymb,mathrsfs,mathpazo}
\usepackage{geometry}
\geometry{a4paper,margin=2cm}
\usepackage{graphicx,array,tabularx,tabularray,colortbl,multirow,multicol,booktabs,enumitem,pifont,ulem,fancybox,diagbox,cancel}
\usepackage{pstricks,pstricks-add,pst-plot,pst-tree,pst-node,pst-text,pst-func,pst-math,pst-bspline,pst-3dplot,pst-eucl,pst-solides3d}
\let\pscalculate\undefined
\usepackage{pst-bezier}
\usepackage{tikz,pgfplots,tkz-tab}
\usetikzlibrary{arrows,arrows.meta,patterns,calc,positioning,trees,decorations.markings,shapes.geometric,shapes.misc,shapes.arrows,intersections,decorations.pathmorphing,decorations.pathreplacing,backgrounds,shadows,babel}
\pgfplotsset{compat=1.18}
\usepgfplotslibrary{fillbetween}
\usepackage[np]{numprint}
\usepackage{listings,makecell,pas-tableur,twoopt}
\let\barre\undefined
\usepackage[pstricks]{tablvar}
\definecolor{darkgreen}{RGB}{0,100,0}
\providecommand{\red}{\color{red}}
\providecommand{\blue}{\color{blue}}
\providecommand{\green}{\color{green}}
\providecommand{\black}{\color{black}}
\providecommand{\white}{\color{white}}
\providecommand{\gray}{\color{gray}}
\providecommand{\cyan}{\color{cyan}}
\providecommand{\magenta}{\color{magenta}}
\providecommand{\yellow}{\color{yellow}}
\providecommand{\violet}{\color{violet}}
\providecommand{\orange}{\color{orange}}
\providecommand{\euro}{\texteuro}
\providecommand{\EUR}{\texteuro}
\providecommand{\ieme}{\textsuperscript{e}}
\providecommand{\degres}{\ensuremath{^\circ}}
\providecommand{\degree}{\ensuremath{^\circ}}
\providecommand{\vv}[1]{\overrightarrow{#1}}
\providecommand{\Coord}[1]{#1}
\providecommand{\decosix}{}
\providecommand{\starredbullet}{\ensuremath{\bullet}}
\pagestyle{empty}
\setlength{\parindent}{0pt}
\usepackage{adjustbox}\DeclareUnicodeCharacter{2236}{:}\DeclareUnicodeCharacter{2190}{\ensuremath{\leftarrow}}
\providecommand{\R}{\mathbb{R}}
\providecommand{\N}{\mathbb{N}}
\providecommand{\Z}{\mathbb{Z}}
\providecommand{\Q}{\mathbb{Q}}
\providecommand{\D}{\mathbb{D}}
\providecommand{\C}{\mathbb{C}}
\providecommand{\e}{{\mathrm{e}}}
\providecommand{\vect}[1]{\overrightarrow{#1}}
\providecommand{\vectt}[1]{\overrightarrow{\mathrm{#1}}}
\providecommand{\vvt}[1]{\overrightarrow{\mathrm{#1}}}
\providecommand{\np}[2][]{{#2~#1}}
\providecommand{\vv}[1]{{\overrightarrow{#1}}}
\renewcommand{\d}{{\mathrm{d}}}
\providecommand{\degres}{{^\circ}}
\providecommand{\degre}{{^\circ}}
\providecommand{\blue}{\color{blue}}
\providecommand{\red}{\color{red}}
\providecommand{\enskip}{\quad}
\providecommand{\llbracket}{\lbrack\!\lbrack}
\providecommand{\rrbracket}{\rbrack\!\rbrack}
\providecommand{\euro}{{\text{€}}}
\providecommand{\no}{{\text{n°}}}
\providecommand{\interval}[2]{{[#1;#2]}}
\providecommand{\intervFO}[2]{{[#1;#2[}}
\providecommand{\dotfill}{\ldots}

\providecommand{\ds}{\displaystyle}
\providecommand{\textding}[1]{#1}
\providecommand{\Oij}{\(\mathrm{O};\vec{\imath},\vec{\jmath}\)}
\providecommand{\Oijk}{\(\mathrm{O};\vec{\imath},\vec{\jmath},\vec{k}\)}
\providecommand{\Ouv}{\(\mathrm{O};\vec{u},\vec{v}\)}
\providecommand{\intervFF}[2]{\ensuremath{[#1;#2]}}
\renewcommand{\euro}{\ifmmode\text{\texteuro}\else\texteuro\fi}
\DeclareUnicodeCharacter{00B0}{\ifmmode^{\circ}\else\textdegree\fi}
\DeclareUnicodeCharacter{20AC}{\ifmmode\text{\texteuro}\else\texteuro\fi}
% Énoncés : Éducation nationale. Transcription et corrigés : APMEP, Denis Vergès et contributeurs. Découpage et classement : CoopMaths, Sébastien Lozano, Éric Elter et contributeurs. Adaptation de la présentation : Maths Explorer.
% CC BY-SA : https://coopmaths.fr/www/about/
% https://coopmaths.fr/alea/static/bac/2026/tex/bac_2026_09_sujet2_metropole_1.tex
\begin{document}
\noindent\textit{APMEP / CoopMaths. CC BY-SA. Présentation adaptée par Maths Explorer.}\par\medskip
\subsection*{Partie A}

Une joueuse de basket constate qu'elle réussit 60\,\% de ses lancers francs au cours d'une séance d'entraînement.

En se basant sur ces statistiques, on modélise la situation comme une succession
d'épreuves aléatoires indépendantes.

\medskip

\begin{enumerate}
	\item Cette joueuse effectue cinq lancers consécutifs.
	\begin{enumerate}
		\item Soit $X$ la variable aléatoire donnant le nombre de lancers réussis.

		Justifier que $X$ suit une loi binomiale dont on précisera les paramètres.
		\item Donner une valeur approchée à $10^{-2}$ près de $P(X \geqslant 4)$.

		Interpréter ce résultat dans le contexte de l'exercice.
	\end{enumerate}
	\item On suppose maintenant que cette sportive doit effectuer $n$ lancers, où $n$ est un entier naturel non nul.

	Soit $Y$ la variable aléatoire donnant le nombre de lancers réussis.

	On admet que la variable aléatoire $Y$ suit la loi binomiale de paramètres $n$ et $p = 0,6$.
	\begin{enumerate}
		\item Montrer que $P(Y \geqslant 1)=1 - 0,4^n$.
		\item En déduire le nombre $n$ de lancers à partir duquel la probabilité d'en
		réussir au moins un dépasse \np{0,9999}.
		\item Recopier et compléter le script de la fonction Python suivante afin qu'elle
		permette de retrouver le résultat de la question \textbf{2. b.}
	\end{enumerate}
\end{enumerate}

\begin{lstlisting}[style=python,xleftmargin=0.3\linewidth,xrightmargin=0.3\linewidth,numbers=none]
def basket ():
	p = 0.6
	n = 1
	while ..... :
		n = n + 1
		p = .....
	return ...
\end{lstlisting}

\subsection*{Partie B}

Pour un autre joueur, on remarque que :
\begin{itemize}
	\item lorsqu'il a réussi un lancer franc, il réussit le suivant avec une probabilité de $0,7$.
	\item lorsqu'il a raté un lancer franc, il réussit le suivant avec une probabilité de $0,4$.
\end{itemize}

On considère que la probabilité de réussir son premier lancer est de  0,5

On note $G_1$ l'évènement : \og le joueur réussit son premier lancer franc \fg.

On a donc $P(G_1) = 0,5$.

On note $G_2$ l'évènement : \og le joueur réussit son deuxième lancer franc \fg .

\medskip

\begin{enumerate}
	\item Calculer $P(G_2)$.

	On pourra utiliser un arbre de probabilité.
\end{enumerate}

Soit $n$ un entier naturel non nul.

On note $G_n$ l'évènement \og le joueur réussit le $n$-ième lancer franc \fg{} et $p_n$, la
probabilité de $G_n$.

On définit ainsi une suite $(p_n)$ de réels.

On admet que pour tout $n \geqslant 1,\: p_{n+1} = 0,3p_n + 0,4$.
\begin{enumerate}[resume]
	\item
	\begin{enumerate}
		\item En utilisant un raisonnement par récurrence, démontrer que

		\begin{center}pour tout $n \geqslant 1,\quad p_n \leqslant p_{n+1}$ \end{center}

		\item Quel argument permet d'affirmer que la suite $(p_n)$ est majorée ?
		\item En déduire que la suite $(p_n)$ est convergente.
	\end{enumerate}
	\item On admet que la limite de $(p_n)$ quand $n$ tend vers $+\infty$ est $\dfrac47$.

	Interpréter cette valeur dans le contexte de l'exercice.
\end{enumerate}

\end{document}
